\begin{flushleft}
Krátké shrnutí a cíl sekce: Tato část popisuje praktická asistivní řešení a postupy, které zlepšují přístupnost výuky pro studenty se speciálními vzdělávacími potřebami (přepis mluveného slova, čtení nahlas, popis obrazového obsahu). Uvedené nástroje a scénáře jsou vybrány podle jejich praktické použitelnosti ve výuce a podle zdrojů popisujících reálné nasazení ve školním prostředí \cite{kb_ai_101_tipu_pdf_04e4a115_page_11_1,kb_ai_101_tipu_pdf_04e4a115_page_7_1}.
\end{flushleft}

\bigskip

Klíčové nástroje a jejich pedagogické role:
\begin{itemize}
  \item Režimy čtení nahlas a asistivní čtečky (např. Microsoft Immersive Reader) — zjednodušují čtení tím, že zvýrazňují text, upravují zobrazení a čtou obsah nahlas; jsou dostupné v mnoha školních aplikacích a pomáhají studentům se čtením a porozuměním \cite{kb_ai_101_tipu_pdf_04e4a115_page_11_1}.
  \item Aplikace pro osoby se zrakovým postižením (SeeingAI) — dokáže v reálném čase popsat objekt či scénu a tím pomáhá nevidomým studentům orientovat se v materiálech a prostředí \cite{kb_ai_101_tipu_pdf_04e4a115_page_11_1}.
  \item Přepis řeči (speech‑to‑text) a překlad v reálném čase (např. Microsoft Translator jako „přepisovač“) — umožňují v hodině zobrazit textový přepis mluveného slova pro studenty se sluchovým postižením nebo pro multijazykové publikum \cite{kb_ai_101_tipu_pdf_04e4a115_page_11_1}.
  \item Funkce školních edukačních aplikací pro sledování pokroku čtení — „Pokrok ve čtení“ v rámci školních nástrojů umí sledovat výkonnost čtení a poskytovat personalizovaná cvičení na základě chyb \cite{kb_ai_101_tipu_pdf_04e4a115_page_7_1}.
\end{itemize}

Praktické scénáře a doporučené postupy:
\begin{enumerate}
  \item Živý přepis přednášky (podpůrný scénář): aktivujte přepis v dostupné aplikaci nebo sdílejte odkaz do Microsoft Translator; studenti vidí přepis v reálném čase na svých zařízeních, což usnadňuje porozumění i pozdější revizi \cite{kb_ai_101_tipu_pdf_04e4a115_page_11_1}.
  \item Asistované čtení materiálů: zařaďte režim čtení nahlas (Immersive Reader) při práci s delšími texty; zvýrazněný a čtený text pomáhá studentům s poruchami čtení lépe sledovat obsah a soustředit se \cite{kb_ai_101_tipu_pdf_04e4a115_page_11_1}.
  \item Personalizované čtecí úlohy: použijte nástroje, které vyhodnocují čtení a generují adaptivní cvičení na slabá místa studenta (příklad: Pokrok ve čtení ve školních edicích nástrojů) \cite{kb_ai_101_tipu_pdf_04e4a115_page_7_1}.
\end{enumerate}

Krátký checklist pro rychlé nasazení v kurzu:
\begin{itemize}
  \item Ověřte dostupnost funkcí v univerzitních licencích a jejich zapnutí v používaných aplikacích (např. Teams / školní aplikace) \cite{kb_ai_101_tipu_pdf_04e4a115_page_7_1}. 
  \item Před první hodinou ohlaste studentům účel použití asistivních nástrojů (přístupnost, podpora učení) a získejte případný souhlas pro nahrávání či ukládání přepisů (general background knowledge).
  \item Před hodinou otestujte technické nastavení (mikrofon, sdílení přepisu, přístupová práva) a mějte záložní plán pro případ výpadku technologie (general background knowledge).
  \item Uvažujte o poskytování alternativních materiálů (PDF s vloženým přepisem, audio verze) pro studenty, kteří nástroj nemohou či nechtějí používat (general background knowledge).
\end{itemize}

Krátké příklady aktivit do hodiny:
\begin{itemize}
  \item „Live captions + otázky“: promítněte krátký výklad (5–7 min) s aktivovaným přepisem; požádejte studenty, aby vyznačili klíčová slovesa nebo fráze a následně otevřete diskusi o rozdílech v porozumění mezi zápisem a mluvením \cite{kb_ai_101_tipu_pdf_04e4a115_page_11_1}. (general background knowledge: konkrétní nastavení přepisu závisí na platformě)
  \item „Orientace s SeeingAI“: nechte studenty ve dvojicích popsat objekt a porovnat vlastní popis se zprávou z aplikace SeeingAI; aktivita ukáže limity a silné stránky strojového popisu \cite{kb_ai_101_tipu_pdf_04e4a115_page_11_1}.
  \item „Krátké čtecí bloky s adaptací“: rozdělte text na krátké části, použijte čtecí režim pro každý blok a po každém bloku dejte cvičení zaměřené na slovní zásobu určenou pro individualizaci (podpůrné cvičení dle nástrojů pro pokrok ve čtení) \cite{kb_ai_101_tipu_pdf_04e4a115_page_7_1}.
\end{itemize}

Poznámky o ochraně osobních údajů a etice (stručně):
\begin{itemize}
  \item Pokud nahráváte nebo ukládáte přepisy či audiozáznamy obsahující osobní údaje, postupujte podle interních pravidel VUT a GDPR‑orientovaných postupů; specifikujte dobu uchování a přístupová práva (general background knowledge).
  \item Transparentně informujte studenty o účelu a rozsahu použití asistivních technologií a nabídněte alternativu pro ty, kteří nástroj nechtějí nebo nemohou využívat (general background knowledge).
\end{itemize}

Doporučení na závěr: začněte s jedním jednoduchým pilotem (např. zapnutí přepisu u jedné přednášky nebo použití čtecího režimu u jednoho textu), zaznamenejte zkušenosti studentů a upravte postupy podle zpětné vazby; nástroje popsané výše jsou především doplňkem pro zlepšení přístupnosti a měly by být integrovány do širší podpory studentů \cite{kb_ai_101_tipu_pdf_04e4a115_page_11_1,kb_ai_101_tipu_pdf_04e4a115_page_7_1}.